\sectionstart
\section{购物链路与外部服务隔离}
\subsection{模拟购物与真实跑腿分开记账}
Mock 链路覆盖合成商品搜索、预算核对、逐项确认、模拟扣款与模拟履约，状态如 \code{mock_paid} 只代表沙箱记录。官方连接器则通过用户授权与美团组件，读取真实地址、询价、准备帮买订单，确认后调用供应商，再由用户进入美团收银台完成真实支付。通用商品搜索仍使用 Mock 数据，不能据跑腿帮买宣称真实餐厅菜单外卖已打通。

\begin{table}[htbp]\centering\small
\caption{两种购物链路的事实边界}
\begin{tabularx}{\textwidth}{@{}L{2.15cm}YY@{}}\toprule
比较项 & Mock 购物 & 美团官方跑腿连接器 \\\midrule
信息来源 & 合成商品与商户 & 授权地址、询价及供应商响应 \\
确认后动作 & 沙箱订单与合成扣款 & 真实帮买订单提交及支付跳转 \\
支付事实 & 仅模拟支付 & 由用户在美团侧完成，须查询真实状态 \\
银行账本 & 写入合成交易与回执 & 不扣合成银行余额 \\
失败处理 & 本地事务与原订单查询 & 保留未知状态、按原凭证核对，不盲目重提 \\\bottomrule
\end{tabularx}
\end{table}

\subsection{用户授权与提交控制}
官方链路须具备 Node 程序、Skill、授权助手和状态目录。每个用户先建立供应商授权，预览绑定用户、地址与询价条件。确认时检查对象归属、令牌、账户版本、过期和预算；较高配送费需要额外确认。订单创建只说明供应商接受了请求，不代表已经支付、接单或履约。

提交超时或结果不明时保留 \code{unknown}，限制同一询价直接重提。调用方应查询原订单状态并解释不确定性，不能因为没有及时收到响应便假定未下单。跨系统幂等能力取决于供应商契约，本地请求去重不能替代对方的执行保证。

真实配送必须向供应商提供收件姓名、电话、地址或坐标，以及物品要求、预算和备注。这是功能所必需的数据传输，与模型汇总字段裁剪是两种不同的数据边界。浏览器展示脱敏地址，不意味着供应商提交报文也只含脱敏信息。

\subsection{组件调用与凭据空间}
连接器通过参数数组启动 Node 子进程，不经 shell 拼接命令。授权缓存目录使用用户 ID 的 SHA-256 派生路径，子进程 HOME、临时目录与状态缓存分开配置；Linux 路径权限设置限制本机普通访问。该设计隔离凭据命名空间，不能等同于独立容器或操作系统级沙箱。

访问 Token 可能由官方授权助手缓存到文件，并在调用时进入子进程环境；部分询价凭证和供应商订单对象会保存到数据库。因此应区分模型密钥、供应商授权 Token、业务订单凭证和公开状态，不使用“所有第三方凭据都不落盘”这样的概括。

地址接口可能出现进程退出码为零但业务失败的情况，系统还需检查供应商业务响应。临时手动地址提供有限兜底，带 HMAC 签名的编码凭证用于完整性验证；签名不等于加密，不能把其中可解码内容当成保密存储。移动端提供打开授权页、手动刷新、重试与重新授权入口。
\source{backend/shopping.py; backend/meituan_paotui.py; scripts/package_meituan_connector.py}

\subsection{邮件服务的真实副作用}
配置 SMTP 时，邮箱验证与密码重置会向真实邮箱发送链接；未配置时，邮件落在权限受限的本机 outbox，属于开发行为，不表示用户已收到邮件。发送过程会向邮件服务提供收件地址。传输是否启用 STARTTLS 取决于配置，不能把可关闭的选项描述为任何环境下都已强制加密。
\source{backend/mailer.py}
